\documentclass[11pt]{article}
\usepackage{fontspec,amsmath,amssymb,amsthm,hyperref}
\usepackage[margin=1in]{geometry}
\setlength{\emergencystretch}{3em}
\title{The conditional one-step KL contraction with bias 6}
\author{LangevinLean mathematical exposition}
\date{}
\begin{document}\maketitle
\noindent Source module: \texttt{UnadjustedLangevin/LSI/OneStep.lean}. Compile this proof document with XeLaTeX or LuaLaTeX. The accompanying HTML places the complete Lean declarations beside every proof.


Assume the displayed LSI, gradient-moment and interpolation differential bounds. For \(0<h\le\alpha/(4L^2)\), a scalar integrating-factor argument proves \(H_{k+1}\le e^{-\alpha h}H_k+6h^2dL^2\).


\section{Absorb a small initial-entropy forcing into the exponential gap}


For \(0\le x\le1/4\), the exponential tangent inequality gives \(e^{-3x/2}\ge1/2\) and \(e^{x/2}\ge1+x/2\). Multiplying the latter by \(e^{-3x/2}\) shows \(e^{-x}-e^{-3x/2}\ge xe^{-3x/2}/2\ge x/4\). Hence

\[e^{-3x/2}+x/4\le e^{-x}.\]

\noindent Lean declarations: \texttt{exp\_three\_halves\_gap}.

\section{Uniform forcing, affine comparison and the coefficient 6}


Assume \(0<\alpha\le L\), \(n\ge0\), \(0<h\le\alpha/(4L^2)\), \(H(0)\ge0\), continuity on \([0,h]\), and right derivatives satisfying


\[2\alpha H(t)\le I(t),\qquad H\prime(t)\le-\tfrac34I(t)+t^2L^2G_0+2tnL^2,\quad0\le t<h,\]
\[G_0\le A:=\frac{4L^2}{\alpha}H(0)+2nL.\]


The step regime implies \(hL\le1/4\) and \(\alpha h\le1/4\). Moreover


\[h^2L^2\frac{4L^2}{\alpha}\le\frac\alpha4,\qquad h^2L^2(2nL)\le\tfrac12hnL^2.\]


Thus the nonnegative upper bound \(A\) gives \(t^2L^2G_0+2tnL^2\le b\), where \(b=(\alpha/4)H(0)+(5/2)hnL^2\ge0\). Substitution of LSI yields \(H\prime\le-(3\alpha/2)H+b\). The affine comparison and \((1-e^{-ah})/a\le h\) give


\[H(h)\le\left(e^{-3\alpha h/2}+\frac{\alpha h}{4}\right)H(0)+\frac52h^2nL^2.\]


The exponential-gap lemma at \(x=\alpha h\), multiplied by \(H(0)\ge0\), absorbs the extra initial-entropy coefficient. Enlarge \(5/2\) to \(6\), using \(h^2nL^2\ge0\), to obtain

\[H(h)\le e^{-\alpha h}H(0)+6h^2nL^2.\]


This scalar argument assumes the displayed production and gradient-moment inequalities; it does not derive them for the Gaussian interpolation.



\noindent Lean declarations: \texttt{one\_step\_entropy\_of\_interpolation}.

\end{document}
